% ======================================================================
%  beamer-keynote.tex -- START HERE. A defense-style academic deck;
%  the theme (beamerthemeAUkeynote.sty) does everything else.
%
%  Anatomy: \section{...} gives you a keynote-style break page for
%  free. Inside frames use normal beamer: itemize, block, alertblock,
%  columns... plus \bignum, \heroimage, \bigequation for set pieces.
%
%  COMPILE WITH LuaLaTeX:   latexmk    (Overleaf: Menu -> Compiler -> LuaLaTeX)
% ======================================================================
\documentclass[aspectratio=169]{beamer}
\usetheme{AUkeynote}

\title[Calabi–Yau Manifolds]{Calabi–Yau Manifolds}
\subtitle{What they are, why they exist, and why they matter}
\author[F. Wang]{Franklin Wang\, · \,Auburn University}
\date{Ph.D. defense · July 2026}

\begin{document}

\CYtitleframe

\begin{frame}{Outline}
  \vspace{0.4em}
  \tableofcontents
\end{frame}

% ======================================================================
\section{Background: shapes and curvature}
% ======================================================================

\begin{frame}{Manifolds: smooth shapes of any dimension}
  \begin{block}{Definition — manifold}
    A space that looks \emph{flat} up close: around every point there are
    coordinates, the way the Earth looks flat from where you stand.
  \end{block}
  \begin{itemize}
    \item The number of coordinates is the \textbf{dimension}; a circle,
          the skin of a doughnut, and spacetime are all manifolds.
    \item We work with \textbf{compact} manifolds: finite and edgeless,
          closed up like a sphere rather than sprawling like a plain.
    \item \textbf{Complex} manifolds run on coordinates $x + iy$ —
          each complex dimension carries two real ones, so a threefold
          $(n = 3)$ is a six-dimensional world.
  \end{itemize}
\end{frame}

\begin{frame}{Curvature is intrinsic: draw a triangle}
  \heroimage[0.42\paperheight]{figures/curvature-trio.png}
  \vspace{0.5em}
  \begin{itemize}
    \item On a \textbf{sphere} the angles sum to \emph{more} than $180^\circ$;
          on a \textbf{saddle}, to \emph{less}; only a \textbf{flat} sheet
          gives exactly $180^\circ$.
    \item An ant that never leaves the surface could discover its world is
          curved — curvature is \textbf{intrinsic}, and varies from point to point.
  \end{itemize}
\end{frame}

\begin{frame}{Ricci curvature: the piece that controls volume}
  \begin{block}{What $\operatorname{Ric}$ measures}
    Sprinkle a small ball of dust and let each speck drift along a
    geodesic: on a sphere the ball \emph{shrinks}, on a saddle it
    \emph{swells}; when $\operatorname{Ric} = 0$ it keeps its volume exactly.
  \end{block}
  \begin{itemize}
    \item A $\operatorname{Ric} = 0$ space is genuinely curved, yet
          volume-balanced.
    \item This is precisely Einstein's equation for empty spacetime:
          \textbf{Ricci-flat spaces are gravity's vacuums}.
  \end{itemize}
\end{frame}

% ======================================================================
\section{Calabi–Yau manifolds and Yau's theorem}
% ======================================================================

\begin{frame}{The definition, and one cheap invariant}
  \begin{block}{Definition — Calabi–Yau manifold}
    A compact complex manifold that is (1) \textbf{Kähler} — metric and
    complex coordinates fit together as smoothly as possible — and
    (2) carries a \textbf{holomorphic volume form} $\Omega$ that vanishes
    nowhere (a trivial canonical bundle, $K_X \cong \mathcal{O}_X$).
  \end{block}
  \begin{itemize}
    \item Feature (2) forces a purely \textbf{topological} number to vanish:
          the first Chern class, $c_1(X) = 0$.
    \item $c_1$ is a whole-number invariant — computable \emph{without
          measuring anything}.
    \item Calabi's question (1954): does that cheap test force a
          Ricci-flat metric to exist?
  \end{itemize}
\end{frame}

\begin{frame}{Yau's theorem}
  \begin{alertblock}{Theorem (Yau 1978 — the Calabi conjecture)}
    Let $(X,\omega)$ be a compact Kähler manifold with $c_1(X) = 0$ in
    $H^2(X,\mathbb{R})$. Then $X$ carries a \emph{unique} Ricci-flat
    Kähler metric $\tilde\omega$ with $[\tilde\omega] = [\omega]$.
  \end{alertblock}
  \vspace{0.4em}
  \bigequation{\operatorname{Ric}(\tilde\omega) \;=\; 0}
  \vspace{0.2em}
  \begin{itemize}
    \item One topological test, $c_1 = 0$, buys the richest geometry of all —
          \textbf{guaranteed}, one metric in every Kähler class.
  \end{itemize}
\end{frame}

\begin{frame}{How it is proved: one nonlinear equation}
  \begin{itemize}
    \item Fix an overall scale: every candidate metric is
          $\omega + i\,\partial\bar\partial\varphi$ for a single unknown
          function $\varphi$.
    \item Ricci-flatness becomes the \textbf{complex Monge–Ampère equation}
          \[ \bigl(\omega + i\,\partial\bar\partial\varphi\bigr)^{n}
             \;=\; e^{F}\,\omega^{n}, \]
          whose unknown sits inside a determinant of its own second derivatives.
    \item Yau's \textbf{a priori estimates} close the continuity method:
          uniform control of $\varphi$, then of its second derivatives,
          forces a solution to exist.
  \end{itemize}
\end{frame}

% ======================================================================
\section{Examples and moduli}
% ======================================================================

\begin{frame}{The simplest one: a torus}
  \begin{columns}[c]
    \begin{column}{0.52\textwidth}
      \begin{itemize}
        \item In one complex dimension the definition selects exactly one
              shape: the torus $\mathbb{C}/\Lambda$.
        \item Its volume form $dz$ descends from the plane; its Ricci-flat
              metric is the flat geometry a sheet of paper keeps when rolled
              up in two directions.
        \item Everything Yau's theorem promises — visible by hand.
      \end{itemize}
    \end{column}
    \begin{column}{0.42\textwidth}
      \heroimage[0.42\paperheight]{figures/cy-torus.png}
    \end{column}
  \end{columns}
\end{frame}

\begin{frame}{The family album: one equation each}
  A degree-$d$ hypersurface in $\mathbb{P}^{n+1}$ is Calabi–Yau exactly
  when $d = n + 2$:
  \vspace{0.5em}
  \begin{center}
    \begin{tabular}{c l l}
      \toprule
      \textbf{dim} & \textbf{Calabi–Yau} & \textbf{why $K_X \cong \mathcal{O}_X$} \\
      \midrule
      $1$ & elliptic curve $\mathbb{C}/\Lambda$ & $dz$ descends \\
      $2$ & K3 surface, e.g.\ $X_4 \subset \mathbb{P}^3$ & $4 = 2 + 2$ \\
      $3$ & quintic $X_5 \subset \mathbb{P}^4$ & $5 = 3 + 2$ \\
      \bottomrule
    \end{tabular}
  \end{center}
  \vspace{0.5em}
  \begin{itemize}
    \item The quintic alone moves in a \textbf{101-parameter family} of
          Ricci-flat worlds.
  \end{itemize}
\end{frame}

\begin{frame}{Seeing a threefold: a slice of the quintic}
  \begin{columns}[c]
    \begin{column}{0.46\textwidth}
      \begin{itemize}
        \item Six real dimensions — too many to draw, not too many to
              \emph{slice}.
        \item Shown: the complex curve $z_1^5 + z_2^5 = 1$ inside the Fermat
              quintic $z_0^5 + \cdots + z_4^5 = 0$ in $\mathbb{P}^4$,
              projected to $3$-space.
        \item Its $25$ patches carry the standard Hanson parametrization.
      \end{itemize}
    \end{column}
    \begin{column}{0.48\textwidth}
      \heroimage[0.62\paperheight]{figures/calabi-yau-quintic.png}
    \end{column}
  \end{columns}
\end{frame}

% ======================================================================
\section{Why it matters}
% ======================================================================

\begin{frame}{String theory: six hidden dimensions}
  \begin{itemize}
    \item String theory needs spacetime to have \textbf{ten} dimensions —
          four we see, six curled up too small to notice.
    \item Candelas–Horowitz–Strominger–Witten (1985): when the hidden six
          form a \textbf{Calabi–Yau threefold}, the four-dimensional physics
          keeps supersymmetry.
    \item Yau's Ricci-flat metric is the vacuum: if string theory is right,
          a shape like the quintic sits at every point of space.
  \end{itemize}
\end{frame}

\begin{frame}{Mirror symmetry: a shape and its mirror}
  \begin{itemize}
    \item Calabi–Yau threefolds come in pairs: each $X$ has a
          \textbf{mirror} $X^{\vee}$ with its two ``shape numbers''
          swapped — ways to \emph{resize} versus ways to \emph{deform}.
    \item String theory cannot tell the pair apart.
    \item Candelas–de la Ossa–Green–Parkes (1991) turned this into
          enumerative geometry: counting curves on $X$ by an integral
          on $X^{\vee}$.
  \end{itemize}
  \vspace{0.4em}
  \begin{center}
    {\KeynoteSemibold\fontsize{20}{24}\selectfont\color{AUorange}
     $2875$ lines\quad·\quad$609{,}250$ conics}\par
    \vspace{0.25em}
    {\KeynoteLight\fontsize{11}{14}\selectfont\color{Quiet}
     on the quintic — counted on its mirror\par}
  \end{center}
\end{frame}

% ======================================================================
\section{Conclusion}
% ======================================================================

\begin{frame}{What to take home}
  \begin{itemize}
    \item A Calabi–Yau manifold passes \textbf{one topological test},
          $c_1 = 0$ — and Yau's theorem rewards it with perfect geometry,
          $\operatorname{Ric} = 0$.
    \item The reward is \emph{guaranteed and unique} in every Kähler class,
          via one nonlinear PDE.
    \item The same shapes carry string compactifications and mirror
          symmetry: \textbf{analysis, algebra, and physics meet here}.
  \end{itemize}
\end{frame}

\CYclosingframe{https://www.linkedin.com/in/franklin-wang-a14b7b327}%
  {Connect on LinkedIn}{zzw0109@auburn.edu}

\end{document}
